\thispagestyle{plain}
\begin{center}
    \Large
    \textbf{Optimized Translation of MATLAB to RISC-V Architecture for Brain-Computer Interfaces}
    
    \vspace{0.4cm}
    \large
    % Thesis Subtitle
    
    \vspace{0.4cm}
    \textbf{Kenan Erol}
    
    \vspace{0.9cm}
    \textbf{Abstract}
\end{center}


Brain computer interfaces (BCIs) offer significant clinical and therapeutic potential for neurological disorders such as epilepsy and stroke. Fully implantable BCIs impose severe limitations on power dissipation to prevent irreversible thermal damage to surrounding brain tissue. As a result, BCI hardware needs to balance real-time signal processing with ultra-low power consumption. The Hardware Architecture for LOw Power BCIs (HALO) project addresses these constraints through a design that combines application-specific hardware accelerators with a low-power RISC-V microcontroller equipped with the RISC-V Vector extension. Despite the powerful capabilities of the HALO architecture, there is a significant toolchain gap between the computational environments preferred by neuroscientists and the low-level constraints of the HALO hardware that proves a hindrance to the widespread adoption of the technology. Algorithm development typically occurs in dynamically typed languages such as MATLAB, whereas HALO hardware requires statically typed, optimized, and vectorized C code. This report details the development of an end-to-end compilation pipeline engineered to bridge this divide. The pipeline uses a TypeScript-based compiler and an RVV-optimized C matrix library to translate MATLAB algorithms into power-efficient RISC-V binaries. The report discusses the architectural design of the pipeline, benchmarks using the Spike instruction-set architecture simulator, investigates execution-cycle metrics, discusses translation benchmarks, and presents next steps for the research.

% bring up to 250 words


